% tikz-tensors-connect.tex -- contraction: which indices a network has, and
% where each one runs.
%
% A tensor's indices are ports. A layout records where each tensor is (its
% stack, the layers k..kb and the sites a..b it covers) and its type says the
% rest (tikz-tensors-nodes.tex); from those, a tensor has
%   up, down     one port per site it covers -- one in the middle, on the side
%                its type names as single
%   left, right  one port per layer it covers
% and a port is a point just inside the tensor, where an index starts, and the
% border it crosses on its way out. Every index here -- one the lattice
% implies (\tnconnect), an open one (\tnopen), one asked for between any two
% ports (\tnjoin) -- runs from a port, so where an index meets a tensor is said
% in one place, for every type, whatever drew it.

% ---- ports -------------------------------------------------------------------
% A height inside tensor #1 at every site it covers, as \tn@ay: its middle,
% unless it is a triangle pointing up or down, which is only as tall as its
% middle near its apex -- then just inside its base.
\def\tn@inner#1{%
  \edef\tn@one{\tn@get{n@#1@one}}%
  \ifx\tn@one\tn@sup \tn@anc{#1}{south}\tn@len\tn@ay{\tn@ay+\tn@inset}%
  \else\ifx\tn@one\tn@sdown \tn@anc{#1}{north}\tn@len\tn@ay{\tn@ay-\tn@inset}%
  \else \tn@anc{#1}{center}\fi\fi}
% Where an index at x = \tn@cx meets tensor #1 of a tree, as \tn@xc: straight
% on if its base reaches that far, else as far out as its base comfortably
% takes an index -- three fifths of the way to its edge. A triangle is half as
% wide at mid-height, where `east' is, as at its base.
\def\tn@meet#1{%
  \tn@anc{#1}{center}\let\tn@mx\tn@ax
  \tn@anc{#1}{east}\tn@len\tn@hw{\tn@ax-\tn@mx}%
  \edef\tn@one{\tn@get{n@#1@one}}%
  \ifx\tn@one\@empty\else \tn@len\tn@hw{2*\tn@hw}\fi
  \tn@len\tn@off{abs(\tn@cx-\tn@mx)}%
  \ifdim\tn@off>\tn@hw
    \tn@len\tn@xc{min(max(\tn@cx,\tn@mx-0.6*\tn@hw),\tn@mx+0.6*\tn@hw)}%
  \else \let\tn@xc\tn@cx \fi}
% \tn@vport{<tensor>}{up|down}: its port on that side in column \tn@cx, as
% \tn@px, \tn@py (where the index starts) and \tn@pe (the y of the border it
% leaves by). In a tree the port is where the tensor's base can take it.
\def\tn@vport#1#2{%
  \let\tn@px\tn@cx
  \ifnum\tn@get{\tn@get{n@#1@s}@tree}=1 \tn@meet{#1}\let\tn@px\tn@xc \fi
  \tn@inner{#1}\let\tn@py\tn@ay
  \edef\tn@t{#2}%
  \ifx\tn@t\tn@sdown \tn@anc{#1}{south}\else \tn@anc{#1}{north}\fi
  \let\tn@pe\tn@ay}
% \tn@hport{<tensor>}{left|right}: its port on that side, as \tn@px (where the
% index starts) and \tn@pe (the x of the border); the height is the layer's.
\def\tn@hport#1#2{%
  \edef\tn@t{#2}%
  \ifx\tn@t\tn@sright
    \tn@anc{#1}{east}\let\tn@pe\tn@ax \tn@len\tn@px{\tn@ax-\tn@inset}%
  \else
    \tn@anc{#1}{west}\let\tn@pe\tn@ax \tn@len\tn@px{\tn@ax+\tn@inset}%
  \fi}

% \tnconnect[along=<style>, down=<style>, apart={<layer>, ...}]{<stack>}
% Every index the grid implies, and nothing else: along a layer between
% neighbouring slots that are both filled (or `-'), border to border; down a
% site from its first tensor to its last, through the ones between, which cover
% it; from an environment into each layer -- to the tensor there, or, if the
% slot is empty, to the border that tensor would have, which is the coordinate
% <stack>-left-<layer> or <stack>-right-<layer>; and up to the dots of a chain
% that goes on. An index that leaves the grid anywhere else is open, and is
% \tnopen's. along= styles the indices along the layers, down= the ones down
% the sites, so that a canonical form's bonds carry arrows and its physical
% indices do not; both are tn edge unless said. apart= names layers whose
% tensors sit side by side without an index between them -- a product, or
% the basis functions of different electrons.
\pgfkeys{/tn/connect/.is family, /tn/connect/along/.initial=tn edge,
         /tn/connect/down/.initial=tn edge, /tn/connect/apart/.initial=}
\newcommand{\tnconnect}[2][]{%
  \pgfkeys{/tn/connect/.cd, along=tn edge, down=tn edge, apart=, #1}%
  \pgfkeysgetvalue{/tn/connect/apart}\tn@apart
  \foreach \tn@l in \tn@apart{\tn@set{#2@apart@\tn@l}{1}}%
  \pgfkeysgetvalue{/tn/connect/along}\tn@along
  \pgfkeysgetvalue{/tn/connect/down}\tn@down
  \edef\tn@ls{\tn@get{#2@layers}}%
  \pgfmathtruncatemacro\tn@nn{\tn@get{#2@n}}%
  \foreach \tn@l in \tn@ls{%
    \tn@layy{#2}{\tn@get{#2@k@\tn@l}}\let\tn@y\tn@ly
    \ifcsname tn@#2@apart@\tn@l\endcsname\else
    \foreach \tn@i in {0,...,\tn@nn}{%
      \pgfmathtruncatemacro\tn@j{\tn@i+1}%
      \tn@occ{#2}{\tn@l}{\tn@i}\let\tn@oa\tn@o \let\tn@ka\tn@kd
      \tn@occ{#2}{\tn@l}{\tn@j}\let\tn@ob\tn@o \let\tn@kb\tn@kd
      \tn@colx{#2}{\tn@i}\let\tn@xa\tn@cx
      \tn@colx{#2}{\tn@j}\let\tn@xb\tn@cx
      \tn@hlink{\tn@along}{#2}%
    }\fi
  }%
  \ifnum\tn@nn>0
  \foreach \tn@i in {1,...,\tn@nn}{%
    \gdef\tn@prev{}%
    \foreach \tn@l in \tn@ls{%
      \tn@occ{#2}{\tn@l}{\tn@i}%
      \ifnum\tn@kd=1
        \ifx\tn@o\tn@prev\else
          \ifx\tn@prev\@empty\else \tn@vlink{#2}{\tn@prev}{\tn@o}{\tn@i}\fi
          \xdef\tn@prev{\tn@o}%
        \fi
      \fi}%
  }%
  \fi
}
% The index between a tensor and the next one down its site. Two that cover
% the same sites meet once on each of them: a gate on a gate, a site on a site.
% Otherwise they meet once, in the middle of the sites they share -- a site
% under a gate, a branch under the node of a tree, an isometry of a MERA under
% the disentangler that overlaps it by half.
\def\tn@vlink#1#2#3#4{% stack, upper, lower, site
  \edef\tn@ua{\tn@get{n@#2@a}}\edef\tn@ub{\tn@get{n@#2@b}}%
  \edef\tn@na{\tn@get{n@#3@a}}\edef\tn@nb{\tn@get{n@#3@b}}%
  \def\tn@draw{0}%
  \ifnum\tn@ua=\tn@na \ifnum\tn@ub=\tn@nb \tn@colx{#1}{#4}\def\tn@draw{1}\fi\fi
  \ifnum\tn@draw=0
    \pgfmathtruncatemacro\tn@oa{max(\tn@ua,\tn@na)}%
    \pgfmathtruncatemacro\tn@ob{min(\tn@ub,\tn@nb)}%
    \ifnum#4=\tn@oa \tn@colx{#1}{(\tn@oa+\tn@ob)/2}\def\tn@draw{1}\fi
  \fi
  \ifnum\tn@draw=1 \tn@facing{\tn@down}{#2}{#3}\fi}
% From the down port of #2 to the up port of #3, under it, both in column
% \tn@cx: straight where the two ports line up, else up, across half way
% between the two borders, and up again.
\def\tn@facing#1#2#3{%
  \tn@vport{#2}{down}\let\tn@xu\tn@px \let\tn@yu\tn@py \let\tn@eu\tn@pe
  \tn@vport{#3}{up}\let\tn@xl\tn@px \let\tn@yl\tn@py
  \ifdim\tn@xu=\tn@xl
    \ifnum\tn@get{\tn@get{n@#2@s}@tree}=1
      \tn@line{#1}{\tn@xl}{\tn@yl}{\tn@xu}{\tn@yu}%
    \else
      \tn@line{#1}{\tn@xu}{\tn@yu}{\tn@xl}{\tn@yl}%
    \fi
  \else
    \tn@len\tn@ym{(\tn@eu+\tn@pe)/2}%
    \tn@turns{#1}{(\tn@xl,\tn@yl) -- (\tn@xl,\tn@ym)
      -- (\tn@xu,\tn@ym) -- (\tn@xu,\tn@yu)}%
  \fi}
% One pair of neighbouring slots along a layer.
\def\tn@hlink#1#2{%
  \ifx\tn@oa\tn@ob\else
  \def\tn@sa{0}\def\tn@sb{0}%
  \ifcase\tn@ka\or\def\tn@sa{1}\or\def\tn@sa{1}\or\or\or\def\tn@sa{1}\or\def\tn@sa{1}\fi
  \ifcase\tn@kb\or\def\tn@sb{1}\or\def\tn@sb{1}\or\or\or\def\tn@sb{1}\or\def\tn@sb{1}\fi
  \ifnum\tn@sa\tn@sb=11
    % a bond: from the right border of the one to the left border of the other
    \def\tn@pa{}\def\tn@pb{}%
    \ifnum\tn@ka=1
      \tn@hport{\tn@oa}{right}\let\tn@xa\tn@px \edef\tn@pa{\tn@get{n@\tn@oa @pt}}\fi
    \ifnum\tn@kb=1
      \tn@hport{\tn@ob}{left}\let\tn@xb\tn@px \edef\tn@pb{\tn@get{n@\tn@ob @pt}}\fi
    \tn@bond{#1}{\tn@xa}{\tn@xb}{\tn@y}{\tn@pa}{\tn@pb}%
  \else
    \ifnum\tn@ka=5 \ifnum\tn@sb=0 \ifnum\tn@kb=4 \else
      \tn@len\tn@xb{\tn@xb-\tn@slot/2}%
      \tn@line{#1}{\tn@xa}{\tn@y}{\tn@xb}{\tn@y}%
      \coordinate (#2-left-\tn@l) at (\tn@xb,\tn@y);
      \tn@reach{\tn@xb}%
    \fi\fi\fi
    \ifnum\tn@kb=6 \ifnum\tn@sa=0 \ifnum\tn@ka=4 \else
      \tn@len\tn@xa{\tn@xa+\tn@slot/2}%
      \tn@line{#1}{\tn@xb}{\tn@y}{\tn@xa}{\tn@y}%
      \coordinate (#2-right-\tn@l) at (\tn@xa,\tn@y);
    \fi\fi\fi
    \ifnum\tn@ka\tn@kb=14
      \tn@hport{\tn@oa}{right}\let\tn@xa\tn@px \tn@len\tn@xb{\tn@xb-\tn@slot/2}%
      \tn@line{#1}{\tn@xa}{\tn@y}{\tn@xb}{\tn@y}%
    \fi
    \ifnum\tn@ka\tn@kb=41
      \tn@hport{\tn@ob}{left}\let\tn@xb\tn@px \tn@len\tn@xa{\tn@xa+\tn@slot/2}%
      \tn@line{#1}{\tn@xa}{\tn@y}{\tn@xb}{\tn@y}%
    \fi
  \fi\fi}

% \tnopen[<style>]{<direction>}{<item>, ...}
% Open indices, up, down, left or right. An index ends at the border the tensor
% in the next slot would have -- past any `|', where the site's index runs on --
% and shows at least `stub' beyond its own tensor. An item is
%   a tensor   its own index that way, one per site it spans up or down and one
%              per layer it spans left or right. The end is <node>-<direction>,
%              or <node>-<direction>-1, -2, ... when there are several
%   a stack    every site's index from its lowest tensor (down) or highest
%              (up), ending at <stack>-<i>-down / -up; or every layer's index
%              from its last tensor (right) or first (left), ending at
%              <stack>-<layer>-right / -left
% A horizontal index that carries a direction takes it from its triangle, as in
% \tnconnect: an open bond of a canonical form points in.
\newcommand{\tnopen}[3][tn edge]{%
  \edef\tn@dir{#2}%
  \foreach \tn@it in {#3}{%
    \ifcsname tn@\tn@it @n\endcsname \tn@openstack{#1}{\tn@it}%
    \else \tn@opennode{#1}{\tn@it}\fi}}
% one index across, from tensor #2 at height #3, ending at coordinate #4
\def\tn@hopen#1#2#3#4{%
  \edef\tn@s{\tn@get{n@#2@s}}\edef\tn@pt{\tn@get{n@#2@pt}}%
  \tn@hport{#2}{\tn@dir}\let\tn@xa\tn@px
  \ifx\tn@dir\tn@sright
    \tn@len\tn@xc{\tn@pe+\tn@stub}%
    \tn@colx{\tn@s}{\tn@get{n@#2@b}+1}\tn@len\tn@xb{\tn@cx-\tn@slot/2}%
    \ifdim\tn@xc>\tn@xb \let\tn@xb\tn@xc \fi
    \tn@reach{\tn@xb}%
    \ifx\tn@pt\tn@sleft \tn@line{#1}{\tn@xb}{#3}{\tn@xa}{#3}%
    \else \tn@line{#1}{\tn@xa}{#3}{\tn@xb}{#3}\fi
  \else
    \tn@len\tn@xc{\tn@pe-\tn@stub}%
    \tn@colx{\tn@s}{\tn@get{n@#2@a}-1}\tn@len\tn@xb{\tn@cx+\tn@slot/2}%
    \ifdim\tn@xc<\tn@xb \let\tn@xb\tn@xc \fi
    \ifx\tn@pt\tn@sright \tn@line{#1}{\tn@xb}{#3}{\tn@xa}{#3}%
    \else \tn@line{#1}{\tn@xa}{#3}{\tn@xb}{#3}\fi
  \fi
  \coordinate (#4) at (\tn@xb,#3);}
% one index up or down, from tensor #2 on site #3, ending at coordinate #4;
% \tn@vopenx draws it at the column <x> instead, walking site #3. It ends at the
% border of the next slot that is not `|'; in a tree, where the tensor's port
% is not over that slot, it turns a corner half way there.
\def\tn@vopen#1#2#3#4{\tn@vopenx{#1}{#2}{#3}{#3}{#4}}
\def\tn@vopenx#1#2#3#4#5{%
  \edef\tn@s{\tn@get{n@#2@s}}%
  \ifx\tn@dir\tn@sdown
    \tn@walk{\tn@s}{#3}{\tn@get{n@#2@kb}+1}{1}%
    \tn@layy{\tn@s}{\tn@kw}\tn@len\tn@yb{\tn@ly+\tn@slot/2}%
  \else
    \tn@walk{\tn@s}{#3}{\tn@get{n@#2@k}-1}{-1}%
    \tn@layy{\tn@s}{\tn@kw}\tn@len\tn@yb{\tn@ly-\tn@slot/2}%
  \fi
  \tn@colx{\tn@s}{#4}%
  \tn@vport{#2}{\tn@dir}%
  \ifdim\tn@px=\tn@cx
    \tn@line{#1}{\tn@cx}{\tn@py}{\tn@cx}{\tn@yb}%
  \else
    \tn@len\tn@ym{(\tn@pe+\tn@yb)/2}%
    \tn@turns{#1}{(\tn@px,\tn@py) -- (\tn@px,\tn@ym)
      -- (\tn@cx,\tn@ym) -- (\tn@cx,\tn@yb)}%
  \fi
  \coordinate (#5) at (\tn@cx,\tn@yb);}
\def\tn@isv{0}
\def\tn@vdir{\def\tn@isv{1}\ifx\tn@dir\tn@sleft\def\tn@isv{0}\fi
  \ifx\tn@dir\tn@sright\def\tn@isv{0}\fi}
\def\tn@opennode#1#2{%
  \tn@vdir
  \edef\tn@s{\tn@get{n@#2@s}}%
  \ifnum\tn@isv=1
    \edef\tn@a{\tn@get{n@#2@a}}\edef\tn@b{\tn@get{n@#2@b}}%
    \edef\tn@one{\tn@get{n@#2@one}}%
    \ifx\tn@one\tn@dir
      \tn@vopenx{#1}{#2}{\tn@a}{(\tn@a+\tn@b)/2}{#2-\tn@dir}%
    \else
    \foreach \tn@i [count=\tn@m] in {\tn@a,...,\tn@b}{%
      \ifnum\tn@a=\tn@b \tn@vopen{#1}{#2}{\tn@i}{#2-\tn@dir}%
      \else \tn@vopen{#1}{#2}{\tn@i}{#2-\tn@dir-\tn@m}\fi}%
    \fi
  \else
    \edef\tn@ka{\tn@get{n@#2@k}}\edef\tn@kb{\tn@get{n@#2@kb}}%
    \foreach \tn@q [count=\tn@m] in {\tn@ka,...,\tn@kb}{%
      \tn@layy{\tn@s}{\tn@q}%
      \ifnum\tn@ka=\tn@kb \tn@hopen{#1}{#2}{\tn@ly}{#2-\tn@dir}%
      \else \tn@hopen{#1}{#2}{\tn@ly}{#2-\tn@dir-\tn@m}\fi}%
  \fi}
\def\tn@openstack#1#2{%
  \tn@vdir
  \edef\tn@ls{\tn@get{#2@layers}}%
  \pgfmathtruncatemacro\tn@nn{\tn@get{#2@n}}%
  \ifnum\tn@isv=1
    \foreach \tn@i in {1,...,\tn@nn}{%
      \gdef\tn@first{}\gdef\tn@lastn{}%
      \foreach \tn@l in \tn@ls{%
        \tn@occ{#2}{\tn@l}{\tn@i}%
        \ifnum\tn@kd=1
          \ifx\tn@first\@empty \xdef\tn@first{\tn@o}\fi
          \xdef\tn@lastn{\tn@o}%
        \fi}%
      \ifx\tn@dir\tn@sdown \let\tn@pick\tn@lastn \else \let\tn@pick\tn@first \fi
      \ifx\tn@pick\@empty\else
        \edef\tn@one{\tn@get{n@\tn@pick @one}}%
        \ifx\tn@one\tn@dir
          \edef\tn@a{\tn@get{n@\tn@pick @a}}\edef\tn@b{\tn@get{n@\tn@pick @b}}%
          \ifnum\tn@i=\tn@a
            \tn@vopenx{#1}{\tn@pick}{\tn@i}{(\tn@a+\tn@b)/2}{#2-\tn@i-\tn@dir}\fi
        \else
          \tn@vopen{#1}{\tn@pick}{\tn@i}{#2-\tn@i-\tn@dir}%
        \fi\fi}%
  \else
    \foreach \tn@l in \tn@ls{%
      \gdef\tn@first{}\gdef\tn@lastn{}%
      \foreach \tn@i in {1,...,\tn@nn}{%
        \tn@occ{#2}{\tn@l}{\tn@i}%
        \ifnum\tn@kd=1
          \ifx\tn@first\@empty \xdef\tn@first{\tn@o}\fi
          \xdef\tn@lastn{\tn@o}%
        \fi}%
      \ifx\tn@dir\tn@sright \let\tn@pick\tn@lastn \else \let\tn@pick\tn@first \fi
      \ifx\tn@pick\@empty\else
        \tn@layy{#2}{\tn@get{#2@k@\tn@l}}%
        \tn@hopen{#1}{\tn@pick}{\tn@ly}{#2-\tn@l-\tn@dir}\fi}%
  \fi}

% \tnopenswap[<style>]{<tensor>}
% The two indices down from a tensor across two sites, exchanged on the way: a
% swap of two fermion legs, drawn by \tnswap, ending where \tnopen{down} would
% end them (<tensor>-down-1 and -2, by where they end). The crossing needs
% room, which a layer of `|' under the tensor gives it.
\newcommand{\tnopenswap}[2][tn edge]{%
  \edef\tn@s{\tn@get{n@#2@s}}\edef\tn@a{\tn@get{n@#2@a}}%
  \tn@walk{\tn@s}{\tn@a}{\tn@get{n@#2@kb}+1}{1}%
  \tn@layy{\tn@s}{\tn@kw}\tn@len\tn@yb{\tn@ly+\tn@slot/2}%
  \tn@anc{#2}{south}\tn@d=\tn@ay\relax\edef\tn@ys{\the\tn@d}%
  \tn@len\tn@drop{\tn@ys-\tn@yb}%
  \tn@colx{\tn@s}{\tn@a}\let\tn@xa\tn@cx
  \tn@colx{\tn@s}{\tn@a+1}\let\tn@xb\tn@cx
  \tnswap[#1]{(\tn@xa,\tn@ys)}{(\tn@xb,\tn@ys)}{\tn@drop}%
  \coordinate (#2-down-1) at (\tn@xa,\tn@yb);
  \coordinate (#2-down-2) at (\tn@xb,\tn@yb);}
